\subsection{Isotypical Modules}

Let A be a ring and S a simple left A-module. Let D be the opposite ring of the endomorphism ring of S; it is a field. Endowed with the actions of A and D, S is an (A,D)-bimodule.

PROPOSITION 1. --- Let M be an A-module. The following properties are equivalent:

\begin{enumerate}
\def\labelenumi{(\roman{enumi})}
\item
  There exists a set I such that M is isomorphic to \(\mathrm{S}^{(\mathrm{I})}\) .
\item
  The module M is the direct sum of a family of submodules isomorphic to S.
\item
  The module M is the sum of a family of submodules isomorphic to S.
\item
  There exists a left vector space V over the field D such that the A-module M is isomorphic to \(S \otimes_{D} V\) .
\end{enumerate}

The equivalence of (i) and (ii) is immediate, and that of (ii) and (iii) follows from Theorem 1 of VIII, p.~56, applied to the case N = 0. Every left vector space over D is isomorphic to a vector space of the form \(\mathrm{D}_{s}^{(\mathrm{I})}\) , where I is a set (II, §7, No.~1, p.~292, Theorem 1). Since the tensor product commutes with direct sums, (i) is equivalent to (iv).

DEFINITION 2. --- An A-module M is isotypical of type S if it has the equivalent properties of Proposition 1. The module M is called isotypical if there exists a simple A-module T such that M is isotypical of type T.

Every isotypical module is semisimple.

PROPOSITION 2. --- If a module is the sum of isotypical submodules of type S, then it is isotypical of type S. The submodules and the quotient modules of an isotypical module of type S are isotypical of type S.

The first assertion follows from the definitions, the second from Corollary 3 of VIII, p.~56.

Remark. --- Every nonzero isotypical module of type S has a quotient module and a submodule isomorphic to S; consequently, if M and M' are nonzero isotypical A-modules of type S, then the group \(\operatorname{Hom}_{\mathrm{A}}(\mathrm{M},\mathrm{M}^{\prime})\) is not reduced to 0.

PROPOSITION 3. --- a) Let M be an isotypical A-module of type S. The A-linear mapping \(\alpha_{\mathrm{M}}: \mathrm{S} \otimes_{\mathrm{D}} \mathrm{Hom}_{\mathrm{A}}(\mathrm{S}, \mathrm{M}) \to \mathrm{M}\) characterized by \(\alpha_{\mathrm{M}}(s \otimes f) = f(s)\) (VIII, p.~59) is bijective.

\begin{enumerate}
\def\labelenumi{\alph{enumi})}
\setcounter{enumi}{1}
\tightlist
\item
  Let V be a left vector space over the field D. The D-linear mapping \(\beta_{\mathrm{V}}: \mathrm{V} \to \mathrm{Hom}_{\mathrm{A}}(\mathrm{S}, \mathrm{S} \otimes_{\mathrm{D}} \mathrm{V})\) defined by \(\beta_{\mathrm{V}}(v)(s) = s \otimes v\) (VIII, p.~59) is bijective.
\end{enumerate}

Denote the left D-vector space \(\mathrm{Hom}_{\mathrm{A}}(\mathrm{S},\mathrm{M})\) by \(\mathcal{H}(\mathrm{M})\) . The A-module M is, by assumption, the direct sum of a family of submodules isomorphic to S. The A-module S is monogenous; to prove that the mapping \(\alpha_{\mathrm{M}}\) is bijective, it therefore suffices to consider the case \(\mathrm{M} = \mathrm{S}\) (VIII, p.~60, Remark 3). Now, \(\mathcal{H}(\mathrm{S})\) is simply the D-vector space \(\mathrm{D}_s\) , and \(\alpha_{\mathrm{S}}\) is simply the isomorphism \(\iota : \mathrm{S} \otimes_{\mathrm{D}} \mathrm{D}_s \to \mathrm{S}\) defined by \(\iota(s \otimes d) = sd\) .

Likewise, to prove b), it suffices to consider the case \(V = D_{s}\) . Since the mapping \(\alpha_{S}\) is bijective, the mapping \(\beta_{D_{s}} = \beta_{\mathcal{H}(S)}\) is too (VIII, p.~60, Remark 2).

